\documentclass[fleqn,onecolumn,a4paper,12pt,titlepage]{article}
\usepackage[showcrop]{geometry}
\geometry{textheight=25.0cm,textwidth=17.0cm,headheight=3cm,headsep=0.5em}
\usepackage{graphicx}
\usepackage{amssymb}
\usepackage{amsmath}
\usepackage[T1]{fontenc}
\usepackage[polish]{babel}
\usepackage{float}
\usepackage[final]{listingsutf8}
\usepackage{fancyhdr}
\usepackage{tikz}
\usetikzlibrary{positioning,arrows.meta,calc,fit,backgrounds}
\usepackage{caption}
\usepackage{datetime2}
\usepackage{enumitem}
\usepackage{hyperref}

\pagestyle{fancy}
\fancyfoot[L,R]{\tiny{\texttt{\symbol{64}}  built: \today{} at \DTMcurrenttime}}

\definecolor{codebg}{RGB}{245,245,244}
\definecolor{comment}{RGB}{0,128,0}
\definecolor{keyword}{RGB}{0,0,180}
\definecolor{string}{RGB}{163,21,21}
\definecolor{linenumber}{RGB}{120,120,120}

\lstdefinestyle{cppstyle}{
    language=C++,
    backgroundcolor=\color{codebg},
    basicstyle=\ttfamily\small,
    keywordstyle=\color{keyword}\bfseries,
    commentstyle=\color{comment}\itshape,
    stringstyle=\color{string},
    numberstyle=\tiny\color{linenumber},
    numbers=left,
    stepnumber=1,
    numbersep=10pt,
    tabsize=4,
    showspaces=false,
    showstringspaces=false,
    breaklines=true,
    breakatwhitespace=true,
    frame=single,
    rulecolor=\color{black},
    captionpos=b
}
\lstset{style=cppstyle}

\newcommand{\placeholderbox}[2]{%
    \begin{center}
    \fbox{\parbox{0.88\textwidth}{\centering\vspace{2.2em}
    \textit{[Miejsce na: #1]}\\\vspace{0.4em}\footnotesize{#2}\vspace{2.2em}}}
    \end{center}
}

\begin{document}

\begin{titlepage}
    \placeholderbox{logo KIA (\texttt{logo\_kia.png})}{Wstaw logo wydziału/katedry}
    \vspace{2cm}
    \centering
    {\LARGE \textsc{Zaawansowane Programowanie w Języku C++} \par}
    \vspace{2cm}
    {\Large \textsc{Raport z wykonania projektu} \par}
    \vspace{2cm}
    {\textsc{GENCLIMB --- aplikacja GTK4 do generowania tras wspinaczkowych Kilter} \par}
    \vfill
    \textbf{[Imię Nazwisko]} $-$ \textbf{[numer indeksu]} $-$ \textbf{[L00]} \par
    \vspace{2cm}
    {\large {\today} \par}
\end{titlepage}

\section*{Wstęp}

\textbf{Kilter Board} to interaktywna ściana wspinaczkowa z podświetlanymi chwytami (dziurkami).
Każda \emph{trasa} to uporządkowany zbiór chwytów z przypisanymi rolami: start, chwyt pośredni,
szczyt (finish) oraz chwyt na nogę. Niniejszy projekt laboratoryjny implementuje natywną
aplikację C++ z interfejsem graficznym GTK4, która:
\begin{itemize}[leftmargin=1.5em]
    \item wczytuje geometrię dziurek ściany z lokalnej bazy SQLite (\texttt{boardlib\_data/kilter.sqlite3}),
    \item generuje nowe trasy przy użyciu wytrenowanego modelu (proces potomny Pythona),
    \item wizualizuje trasę na planszy,
    \item pozwala użytkownikowi ocenić trasę (1--5) i zapisać ją do lokalnej bazy wyników.
\end{itemize}

Kod aplikacji znajduje się w katalogu \texttt{kilter\_cpp\_gui/}. Projekt został
zorganizowany modułowo: nagłówki w \texttt{include/}, implementacje w \texttt{src/},
konfiguracja budowania w \texttt{CMakeLists.txt}. Uruchomienie wymaga katalogu
głównego repozytorium jako bieżącego katalogu roboczego (ścieżki względne do baz
i skryptów Pythona).

\section*{Baza danych}

Aplikacja korzysta z \textbf{dwóch instancji SQLite} --- osobnej bazy tylko do odczytu
oraz zapisywalnej bazy wyników użytkownika.

\subsection*{Baza planszy (\texttt{kilter.sqlite3})}

Baza \emph{boardlib} zawiera co najmniej trzy tabele wykorzystywane przez zapytanie
w module \texttt{board\_db}:
\begin{itemize}[leftmargin=1.5em]
    \item \texttt{holes} --- identyfikator i współrzędne każdej dziurki na panelu,
    \item \texttt{placements} --- powiązanie dziurki z układem (\texttt{layout\_id}, \texttt{set\_id}),
    \item \texttt{product\_sizes\_layouts\_sets} --- mapowanie wariantu produktu
          (\texttt{product\_size\_id}) na układ ściany.
\end{itemize}
Aplikacja odczytuje dziurki dla wybranego układu (token layoutu, np.\ 1270) poprzez
złączenie (\texttt{JOIN}) tych tabel. Funkcja \texttt{product\_size\_id\_from\_layout\_token}
mapuje token modelu na identyfikator produktu w bazie boardlib.

\subsection*{Baza wyników (\texttt{saved\_routes.db})}

Zapisywalna baza tworzona automatycznie w \texttt{kilter\_cpp\_gui/data/} przechowuje
historię tras zaakceptowanych przez użytkownika. Tabela \texttt{saved\_routes} zawiera
pola: \texttt{id}, \texttt{created\_at}, \texttt{layout\_id}, \texttt{grade\_id},
\texttt{source} (model lub dataset), \texttt{rating}, \texttt{hole\_count},
\texttt{tokens}, \texttt{hole\_ids}.

Operacje na bazie wyników realizuje moduł \texttt{results\_db}:
inicjalizacja schematu (\texttt{init\_results\_db}), zapis (\texttt{save\_route\_to\_db})
oraz odczyt ostatnich wpisów (\texttt{load\_recent\_routes}, sortowanie
\texttt{ORDER BY id DESC}).

\subsection*{Diagram bazy danych}

Na rysunku~\ref{fig:db_diagram} przedstawiono uproszczony schemat relacji między
tabelami obu baz. Strzałki oznaczają klucze obce logiczne wykorzystywane w zapytaniach
aplikacji (niekoniecznie zadeklarowane jako \texttt{FOREIGN KEY} w pliku SQLite).

\begin{figure}[H]
    \centering
    \begin{tikzpicture}[
        table/.style={draw, rounded corners=2pt, align=center, font=\footnotesize\ttfamily,
                      inner sep=6pt, fill=white, text width=3.6cm},
        dbframe/.style={draw, dashed, rounded corners=4pt},
        rel/.style={-{Stealth[length=2mm]}, semithick, shorten >=3pt, shorten <=3pt},
        rellbl/.style={font=\scriptsize, fill=white, inner sep=2pt}
    ]
        \node[table] (psls) at (0, 0) {%
            \textbf{product\_sizes\_layouts\_sets}\\[3pt]
            product\_size\_id\\
            layout\_id, set\_id
        };
        \node[table, right=2.6cm of psls] (placements) {%
            \textbf{placements}\\[3pt]
            layout\_id, set\_id\\
            hole\_id
        };
        \node[table, right=2.6cm of placements] (holes) {%
            \textbf{holes}\\[3pt]
            id\\
            x, y
        };

        \draw[rel] (psls.east) -- (placements.west);
        \draw[rel] (placements.east) -- (holes.west);
        \node[rellbl] at ($(psls)!0.5!(placements) + (0,0.55)$) {layout\_id, set\_id};
        \node[rellbl] at ($(placements)!0.5!(holes) + (0,0.55)$) {hole\_id};

        \node[dbframe, fit=(psls) (holes), inner xsep=16pt, inner ysep=12pt] (boardframe) {};
        \node[font=\small, anchor=south west, yshift=2pt] at (boardframe.north west) {%
            \texttt{boardlib\_data/kilter.sqlite3} (odczyt)};

        \node[table, below=3.2cm of placements] (saved) {%
            \textbf{saved\_routes}\\[3pt]
            id, created\_at\\
            layout\_id, grade\_id\\
            source, rating\\
            hole\_count, tokens
        };
        \node[dbframe, fit=(saved), inner xsep=16pt, inner ysep=12pt] (resframe) {};
        \node[font=\small, anchor=south west, yshift=2pt] at (resframe.north west) {%
            \texttt{kilter\_cpp\_gui/data/saved\_routes.db} (zapis)};
    \end{tikzpicture}
    \caption{Uproszczony diagram tabel baz danych wykorzystywanych przez aplikację.}
    \label{fig:db_diagram}
\end{figure}

\subsection*{Implementacja --- połączenie z bazą}

Inicjalizacja lokalnej bazy wyników oraz odczyt dziurek planszy ilustrują listingi
~\ref{lst:results_db} i~\ref{lst:board_db}. Dostęp do SQLite realizowany jest przez
bibliotekę \texttt{libsqlite3} (CMake: \texttt{PkgConfig::SQLITE3}). Parametryzowane
zapytania (\texttt{sqlite3\_prepare\_v2}, \texttt{sqlite3\_bind\_*}) zapobiegają
wstrzyknięciu SQL.

\lstinputlisting[caption={Inicjalizacja tabeli \texttt{saved\_routes} w lokalnej bazie wyników},
                 label={lst:results_db}]{listings/results_db_init.cpp}

\lstinputlisting[caption={Odczyt pozycji dziurek z bazy boardlib (JOIN trzech tabel)},
                 label={lst:board_db}]{listings/board_db_query.cpp}

\newpage

\section*{Struktura aplikacji}

\subsection*{Organizacja modułów}

Po reorganizacji projektu kod podzielono na moduły z wyraźnym podziałem
\texttt{include/} / \texttt{src/}:

\begin{center}
\small
\begin{tabular}{l|l}
\textbf{Moduł} & \textbf{Odpowiedzialność} \\
\hline
\texttt{types.hpp} & Struktury \texttt{Hole}, \texttt{SavedRoute}, \texttt{AppState} \\
\texttt{board\_db} & Odczyt dziurek z \texttt{kilter.sqlite3} \\
\texttt{results\_db} & Zapis i odczyt historii tras użytkownika \\
\texttt{generator\_process} & Klasa \texttt{GeneratorProcess} --- IPC z workerem Pythona \\
\texttt{route} & Parsowanie tokenów tras i przypisywanie ról chwytów \\
\texttt{ui} & Interfejs GTK4, rysowanie planszy (Cairo), obsługa zdarzeń \\
\texttt{main.cpp} & Inicjalizacja ścieżek, uruchomienie \texttt{GtkApplication} \\
\end{tabular}
\end{center}

Program jest napisany obiektowo w sensie modularnym: logika enkapsulowana w klasach
(np.\ \texttt{GeneratorProcess}) oraz w funkcjach powiązanych z konkretnym obszarem
danych (bazy, trasa, UI). Stan aplikacji zebrany jest w strukturze \texttt{AppState}.

\subsection*{Diagram modułów}

\begin{figure}[H]
    \centering
    \begin{tikzpicture}[
        mod/.style={draw, rounded corners=2pt, minimum width=2.5cm, minimum height=0.9cm,
                    align=center, font=\small, fill=white},
        ext/.style={draw, dashed, rounded corners=2pt, minimum width=2.5cm, minimum height=0.9cm,
                    align=center, font=\small, fill=gray!6},
        arr/.style={-{Stealth[length=2mm]}, semithick},
        arrd/.style={-{Stealth[length=2mm]}, semithick, dashed},
        lbl/.style={font=\scriptsize, fill=white, inner sep=2pt}
    ]
        \node[mod] (main) at (0, 0) {main.cpp};

        \node[mod] (ui) at (-4.2, -2.2) {ui};
        \node[mod] (board) at (0, -2.2) {board\_db};
        \node[mod] (results) at (4.2, -2.2) {results\_db};

        \node[mod] (route) at (-5.8, -4.6) {route};
        \node[mod] (gen) at (-2.6, -4.6) {Generator\\Process};
        \node[ext] (py) at (1.2, -4.6) {Python\\worker};

        \draw[arr] (main.south) -- ++(0,-0.55) -| (ui.north);
        \draw[arr] (main.south) -- ++(0,-0.55) -| (board.north);
        \draw[arr] (main.south) -- ++(0,-0.55) -| (results.north);

        \draw[arr] (ui.south) -- ++(0,-0.45) -| (route.north);
        \draw[arr] (ui.south) -- ++(0,-0.45) -| (gen.north);

        \draw[arr] (ui.east) to[out=0, in=180, looseness=1.1] (board.west);
        \draw[arr] (ui.east) to[out=10, in=170, looseness=1.3] (results.west);

        \draw[arrd] (gen.east) -- node[lbl, above, pos=0.55] {pipe JSON} (py.west);
    \end{tikzpicture}
    \caption{Zależności między modułami aplikacji \texttt{kilter\_gui}.}
    \label{fig:modules}
\end{figure}

\subsection*{Komunikacja z procesem generatora}

Wymaganie poziomu~II (dwa procesy komunikujące się przez potok) realizuje klasa
\texttt{GeneratorProcess}. Zamiast uruchamiać interpreter Pythona przy każdym kliknięciu,
aplikacja:
\begin{enumerate}[leftmargin=1.5em]
    \item tworzy parę potoków (\texttt{pipe}),
    \item wywołuje \texttt{fork()} i w procesie potomnym \texttt{execl()} na skrypt
          \texttt{kilter\_dl/generate\_server.py},
    \item przekierowuje \texttt{stdin}/\texttt{stdout} workera na końce potoków,
    \item wymienia komunikaty JSON linia po linii (żądanie \texttt{generate},
          odpowiedź \texttt{ok}/\texttt{error}, na koniec \texttt{shutdown}).
\end{enumerate}

Listing~\ref{lst:generator_ipc} pokazuje kluczowe fragmenty startu i wysyłania żądania.
Obsługa w GUI znajduje się w \texttt{on\_generate\_clicked} (listing~\ref{lst:ui_ipc}).

\lstinputlisting[caption={Uruchomienie workera Pythona i komunikacja przez potoki (IPC)},
                 label={lst:generator_ipc}]{listings/generator_process_ipc.cpp}

\lstinputlisting[caption={Wysłanie żądania generacji trasy z poziomu GUI},
                 label={lst:ui_ipc}]{listings/ui_generate_request.cpp}

% \subsection*{Spełnienie kryteriów oceny}

% \begin{center}
% \small
% \begin{tabular}{p{3.2cm}|p{10.5cm}}
% \textbf{Poziom} & \textbf{Realizacja w projekcie} \\
% \hline
% I (3.0) &
% Baza SQLite z $\geq 3$ tabelami (boardlib); odczyt dziurek i zapis tras użytkownika;
% program modularny/obiektowy. \emph{Uwaga:} w kodzie brak jawnych statystyk
% (średnia, mediana, odchylenie) oraz interaktywnego sortowania listy wyników
% poza \texttt{ORDER BY id DESC}. \\
% \hline
% II (4.0) &
% Spełnione: (a)~baza jako instancja SQLite, nie plik tekstowy;
% (b)~dwa procesy (GUI C++ + worker Python) komunikujące się przez potoki. \\
% \hline
% III (5.0) &
% Wariant alternatywny: poziom~I + przemyślany interfejs graficzny GTK4
% (plansza Cairo, panele sterowania, dialog potwierdzenia zapisu).
% Dodatkowo spełnione są też punkty IIa i IIb. \\
% \end{tabular}
% \end{center}

\newpage

\section*{Działanie aplikacji}

\subsection*{Budowanie i uruchomienie}

\begin{verbatim}
./kilter_cpp_gui/run.sh          # build + run z katalogu głównego repo
./kilter_cpp_gui/run.sh build    # tylko kompilacja
\end{verbatim}

Aplikacja wymaga: \texttt{boardlib\_data/kilter.sqlite3}, środowiska Python
\texttt{kilter\_dl/.venv}, co najmniej jednego checkpointu w \texttt{kilter\_dl/checkpoints/}
oraz plików map tokenów w \texttt{dataset/}.

\subsection*{Przepływ użytkownika}

\begin{enumerate}[leftmargin=1.5em]
    \item Użytkownik wybiera układ ściany (\textbf{Board}) i trudność (\textbf{Grade}).
    \item Przycisk \emph{Generate Route + Render} uruchamia (lub ponownie wykorzystuje)
          proces generatora i wysyła żądanie JSON z parametrami dekodowania.
    \item Odpowiedź zawiera tokeny trasy; moduł \texttt{route} przypisuje role chwytom
          (start --- zielony, pośredni --- niebieski, finish --- fioletowy, noga --- pomarańczowy).
    \item Plansza po prawej stronie rysuje wszystkie dziurki oraz aktywną trasę (Cairo).
    \item Użytkownik ustawia ocenę 1--5 i klika \emph{Save route + rating to DB};
          po potwierdzeniu w dialogu rekord trafia do \texttt{saved\_routes.db}.
    \item Panel \emph{Recently saved routes} odświeża listę ostatnich 10 zapisów.
    \item Alternatywnie: \emph{Load Random Route from Dataset} pobiera trasę ze zbioru
          bez uruchamiania modelu.
\end{enumerate}

\subsection*{Zrzuty ekranu}

\begin{figure}[H]
    \placeholderbox{zrzut ekranu --- główne okno aplikacji}{Widok po wygenerowaniu trasy:
    panel sterowania po lewej, plansza z kolorowymi chwytami po prawej}
    \caption{Główne okno aplikacji GENCLIMB po wygenerowaniu trasy.}
    \label{fig:main_window}
\end{figure}

\begin{figure}[H]
    \placeholderbox{zrzut ekranu --- dialog zapisu trasy}{Potwierdzenie zapisu z oceną
    użytkownika (1--5) do lokalnej bazy}
    \caption{Dialog potwierdzenia zapisu trasy do bazy wyników.}
    \label{fig:save_dialog}
\end{figure}

\begin{figure}[H]
    \placeholderbox{zrzut ekranu --- panel historii zapisów}{Lista ostatnich tras
    z \texttt{saved\_routes.db} (layout, grade, źródło, liczba chwytów, ocena)}
    \caption{Panel historii zapisanych tras.}
    \label{fig:history}
\end{figure}

\section*{Podsumowanie}

Zrealizowano aplikację C++ z graficznym interfejsem GTK4 do pracy z trasami
wspinaczkowymi na ścianach Kilter. Projekt łączy odczyt danych geometrycznych
z bazy boardlib, generację tras przez długożyjący proces potomny (IPC przez potoki)
oraz trwały zapis ocen użytkownika w osobnej bazie SQLite.

\textbf{Najciekawszy element} --- architektura z procesem workera Pythona:
jednokrotne załadowanie modelu PyTorch skraca czas kolejnych generacji w porównaniu
z uruchamianiem nowego interpretera przy każdym żądaniu. Protokół JSON linia-po-linii
upraszcza debugowanie kanału komunikacji.

\textbf{Najtrudniejszy element} --- spójne mapowanie tokenów modelu na identyfikatory
dziurek i role chwytów oraz synchronizacja wielu ścieżek danych (model, dataset, baza
planszy) przy zachowaniu czytelnej struktury modułowej po refaktoryzacji katalogów
\texttt{include/} i \texttt{src/}.

\textbf{Obszar do rozbudowy} --- zgodnie z kryteriami poziomu~I warto dodać moduł
statystyk (np.\ średnia/mediana ocen i liczby chwytów w \texttt{saved\_routes})
oraz interaktywne sortowanie widocznej historii zapisów.

\end{document}
